% GENERATED FILE. Edit canonical lessons, then run npm run generate:books.
% canonical-source-hash: fnv1a64:de2f79c9b7ad280f
% canonical-lessons: JA-C192-kouban, JA-C192-omawarisan, JA-C192-keisatsu, JA-C192-tazuneru, JA-C192-annai

\chapter{Asking a Police Officer}
\label{ch:ja-asking-a-police-officer}
\hlchaptermodality{192}

\begin{chapteropening}
\textbf{By the end of this chapter:} \emph{I can ask a police officer the way.}

Complete the last lesson of chapter 192: I can ask a police officer the way.
\end{chapteropening}

\section[\texorpdfstring{\ja{こうばん} (kōban)}{こうばん (kōban)}]{\ja{こうばん} (kōban) — police box}

\label{lesson:JA-C192-kouban}



\begin{quote}
Before the new one: say the Japanese for pharmacy, then the Japanese for to find.
\end{quote}

\subsection*{You'll want to know: \ja{こうばん}}
\textbf{\ja{こうばん}} — \emph{kōban} — \textquotedblleft{}police box\textquotedblright{}.

\begin{quote}
\textbf{\ja{えきの} \ja{まえに} \ja{こうばんが} \ja{あります。}} — \emph{eki no mae ni kōban ga arimasu} — \textquotedblleft{}There is a police box in front of the station.\textquotedblright{}
\end{quote}

\begin{cousinweb}
\textbf{\ja{こうばん}} (\emph{kōban}) is Sino-Japanese (kango): two Chinese characters read with their on'yomi, readings borrowed from Chinese. They are \emph{kō}, \textquotedblleft{}mingle, coming and going\textquotedblright{} (Middle Chinese \emph{kaew}), and \emph{ban}, \textquotedblleft{}turn, number in a series\textquotedblright{} (KANJIDIC2; Middle Chinese in Baxter's transcription).
\end{cousinweb}

\begin{grammarlens}[title={What you've built}]
One more word for everyday life at A2.
\end{grammarlens}

\subsection*{Your turn}
\begin{itemize}
\raggedright
  \item \emph{Say it:} \emph{kōban}
  \item \emph{Say it:} \emph{kōban}, once more
  \item \emph{Say it:} \emph{eki no mae ni kōban ga arimasu} — There is a police box in front of the station.
  \item \emph{Recall:} say the Japanese for pharmacy, then the Japanese for to find, then say \emph{kōban} again
\end{itemize}

\begin{tcolorbox}[breakable,colback=teal!4,colframe=teal!35!black,title={Before you move on}]
What does \ja{こうばん} mean? (\textquotedblleft{}Police box\textquotedblright{}.) Say it once more.
\end{tcolorbox}
\section[\texorpdfstring{\ja{おまわりさん} (omawarisan)}{おまわりさん (omawarisan)}]{\ja{おまわりさん} (omawarisan) — police officer}

\label{lesson:JA-C192-omawarisan}



\begin{quote}
Before the new one: say the Japanese for to find, then the Japanese for police box.
\end{quote}

\subsection*{You'll want to know: \ja{おまわりさん}}
\textbf{\ja{おまわりさん}} — \emph{omawarisan} — \textquotedblleft{}police officer\textquotedblright{}.

\begin{quote}
\textbf{\ja{こうばんに} \ja{おまわりさんが} \ja{います。}} — \emph{kōban ni omawarisan ga imasu} — \textquotedblleft{}There is a police officer at the police box.\textquotedblright{}
\end{quote}

\begin{cousinweb}
\textbf{\ja{おまわりさん}} (\emph{omawarisan}) is built from polite \emph{o-}, native \emph{mawari}, \textquotedblleft{}going round, a patrol\textquotedblright{}, from \emph{mawaru}, \textquotedblleft{}to go round\textquotedblright{}, and the polite \emph{-san} (JMdict, KANJIDIC2).
\end{cousinweb}

\begin{grammarlens}[title={What you've built}]
One more word for everyday life at A2.
\end{grammarlens}

\subsection*{Your turn}
\begin{itemize}
\raggedright
  \item \emph{Say it:} \emph{omawarisan}
  \item \emph{Say it:} \emph{omawarisan}, once more
  \item \emph{Say it:} \emph{kōban ni omawarisan ga imasu} — There is a police officer at the police box.
  \item \emph{Recall:} say the Japanese for to find, then the Japanese for police box, then say \emph{omawarisan} again
\end{itemize}

\begin{tcolorbox}[breakable,colback=teal!4,colframe=teal!35!black,title={Before you move on}]
What does \ja{おまわりさん} mean? (\textquotedblleft{}Police officer\textquotedblright{}.) Say it once more.
\end{tcolorbox}
\section[\texorpdfstring{\ja{けいさつ} (keisatsu)}{けいさつ (keisatsu)}]{\ja{けいさつ} (keisatsu) — police}

\label{lesson:JA-C192-keisatsu}



\begin{quote}
Before the new one: say the Japanese for police box, then the Japanese for police officer.
\end{quote}

\subsection*{You'll want to know: \ja{けいさつ}}
\textbf{\ja{けいさつ}} — \emph{keisatsu} — \textquotedblleft{}police\textquotedblright{}.

\begin{quote}
\textbf{\ja{けいさつに} \ja{でんわを} \ja{します。}} — \emph{keisatsu ni denwa o shimasu} — \textquotedblleft{}I phone the police.\textquotedblright{}
\end{quote}

\begin{cousinweb}
\textbf{\ja{けいさつ}} (\emph{keisatsu}) is Sino-Japanese (kango): two Chinese characters read with their on'yomi, readings borrowed from Chinese. They are \emph{kei}, \textquotedblleft{}admonish, commandment\textquotedblright{} (Middle Chinese \emph{kjaengX}) and \emph{satsu}, \textquotedblleft{}judge, understand\textquotedblright{} (\emph{tsrheat}) (KANJIDIC2; Middle Chinese in Baxter's transcription).
\end{cousinweb}

\begin{grammarlens}[title={What you've built}]
One more word for everyday life at A2.
\end{grammarlens}

\subsection*{Your turn}
\begin{itemize}
\raggedright
  \item \emph{Say it:} \emph{keisatsu}
  \item \emph{Say it:} \emph{keisatsu}, once more
  \item \emph{Say it:} \emph{keisatsu ni denwa o shimasu} — I phone the police.
  \item \emph{Recall:} say the Japanese for police box, then the Japanese for police officer, then say \emph{keisatsu} again
\end{itemize}

\begin{tcolorbox}[breakable,colback=teal!4,colframe=teal!35!black,title={Before you move on}]
What does \ja{けいさつ} mean? (\textquotedblleft{}Police\textquotedblright{}.) Say it once more.
\end{tcolorbox}
\section[\texorpdfstring{\ja{たずねる} (tazuneru)}{たずねる (tazuneru)}]{\ja{たずねる} (tazuneru) — to ask; to inquire, to look for}

\label{lesson:JA-C192-tazuneru}



\begin{quote}
Before the new one: say the Japanese for police officer, then the Japanese for police.
\end{quote}

\subsection*{You'll want to know: \ja{たずねる}}
\textbf{\ja{たずねる}} — \emph{tazuneru} — \textquotedblleft{}to ask; to inquire, to look for\textquotedblright{}.

\begin{quote}
\textbf{\ja{おまわりさんに} \ja{みちを} \ja{たずねます。}} — \emph{omawarisan ni michi o tazunemasu} — \textquotedblleft{}I ask a police officer the way.\textquotedblright{}
\end{quote}

Polite \textbf{\ja{たずねます}}, \emph{tazunemasu}; an ichidan verb.

\begin{cousinweb}
\textbf{\ja{たずねる}} (\emph{tazuneru}) is a native Japanese word (wago), inherited rather than borrowed from Chinese; its kanji was attached later as a kun'yomi (JMdict, KANJIDIC2).
\end{cousinweb}

\begin{grammarlens}[title={What you've built}]
One more word for everyday life at A2.
\end{grammarlens}

\subsection*{Your turn}
\begin{itemize}
\raggedright
  \item \emph{Say it:} \emph{tazuneru}
  \item \emph{Say it:} \emph{tazuneru}, once more
  \item \emph{Say it:} \emph{omawarisan ni michi o tazunemasu} — I ask a police officer the way.
  \item \emph{Recall:} say the Japanese for police officer, then the Japanese for police, then say \emph{tazuneru} again
\end{itemize}

\begin{tcolorbox}[breakable,colback=teal!4,colframe=teal!35!black,title={Before you move on}]
What does \ja{たずねる} mean? (\textquotedblleft{}To ask; to inquire, to look for\textquotedblright{}.) Say it once more.
\end{tcolorbox}
\section[\texorpdfstring{\ja{あんない} (annai)}{あんない (annai)}]{\ja{あんない} (annai) — guidance; showing the way, information}

\label{lesson:JA-C192-annai}



\begin{quote}
Before the new one: say the Japanese for police, then the Japanese for to ask.
\end{quote}

\subsection*{You'll want to know: \ja{あんない}}
\textbf{\ja{あんない}} — \emph{annai} — \textquotedblleft{}guidance; showing the way, information\textquotedblright{}.

\begin{quote}
\textbf{\ja{おまわりさんが} \ja{みちを} \ja{あんないします。}} — \emph{omawarisan ga michi o annai shimasu} — \textquotedblleft{}The police officer shows the way.\textquotedblright{}
\end{quote}

\begin{cousinweb}
\textbf{\ja{あんない}} (\emph{annai}) is Sino-Japanese (kango): two Chinese characters read with their on'yomi, readings borrowed from Chinese. They are \emph{an}, \textquotedblleft{}plan, idea\textquotedblright{} (Middle Chinese \emph{'anH}) and \emph{nai}, \textquotedblleft{}inside\textquotedblright{} (\emph{nwojH}) (KANJIDIC2; Middle Chinese in Baxter's transcription).
\end{cousinweb}

\begin{grammarlens}[title={What you've built}]
One more word for everyday life at A2.
\end{grammarlens}

\subsection*{Your turn}
\begin{itemize}
\raggedright
  \item \emph{Say it:} \emph{annai}
  \item \emph{Say it:} \emph{annai}, once more
  \item \emph{Say it:} \emph{omawarisan ga michi o annai shimasu} — The police officer shows the way.
  \item \emph{Recall:} say the Japanese for police, then the Japanese for to ask, then say \emph{annai} again
\end{itemize}

\begin{tcolorbox}[breakable,colback=teal!4,colframe=teal!35!black,title={Before you move on}]
What does \ja{あんない} mean? (\textquotedblleft{}Guidance; showing the way, information\textquotedblright{}.) Say it once more.
\end{tcolorbox}
